% Q12_gold_answers.tex — GRADER CALIBRATION FIXTURE (correct, paraphrased).
% A drop-in replacement for paper/Q12_answers.tex (spec §4.7, §6.12): every answer is
% correct, but the wording, the seeds and the baseline are deliberately far from the
% reference, so the grader must recognise correctness by substance (the rubric's
% probes), not by matching text. B1 uses a baseline of 0.4 and its own seeds; the cards
% word every field their own way. B2 is not graded in this mode.
%
% EXPECTED[B1][hypothesis-precision][4] [budget-parity][4] [variance][4] [contamination][4] [cheat-audit][4] [disclosure][4]: correct. Probe B recomputes every product (4320, 80, 160) and the totals against the registry (4,400 / 620 and 180 / 20); Probe M at the stated inputs prints 0.2854; Probe S finds every seed list disjoint; the markers appear once. The report's first line is PREREG Q12: UNLOCKED — every method dimension >= 3. No finding.
% EXPECTED[B3][card-fidelity][4] [B4][card-fidelity][4] [B5][card-fidelity][4]: every field agrees with its pinned range; Variance and Artifact fields correctly say what each range does not give. No finding.
% Expected per-problem grades: B1 = 4, B3 = 4, B4 = 4, B5 = 4 (3 is tolerated only with a declared-gap note and no flaw).
% A calibration run FAILS if any graded problem scores <= 2, if B1 is HELD, or if any
% finding claims a false product, a wrong MDE, an overlap of seeds or a contradiction
% with a range.
\documentclass[11pt]{article}
\input{preamble}

\title{Q12 --- Practice I-C + Lab 1 \\ \large Written answers}
\author{learner}
\date{}

\begin{document}
\maketitle
\sloppy

% PREREG-BEGIN
\section*{Problem B1 (pre-registration)}\label{ans:B1}
\begin{answer}
\textbf{Claim} Three two-sided hypotheses, each refuted by a paired-difference interval
that contains zero at its Holm level: (H1) repairs from feedback (A2) change the rate
against independent samples (A1); (H2) best-first search (A3) changes it against A1;
(H3) the A2 $-$ A1 effect is different in the two languages.

\textbf{Metric and estimator} A task in a repeat is solved when some candidate within
its 8 calls passes the gate; per arm, the mean over tasks of the fraction of repeats
solved; per comparison, the mean paired difference, with a percentile bootstrap over
tasks (repeats kept within their task) and McNemar on each repeat as a check.

\textbf{Arms and policies} A1 \texttt{IndependentSampler}, A2 \texttt{SequentialRepair},
A3 \texttt{BestFirst}, all from \texttt{harness/q/fvh\_q/policies.py}; prompt hashes
recorded once the file is final. The Claude arm runs A2 on 20 tasks, 10 per language
and 5 per stratum within each language, drawn with seed 7711.

\textbf{Tasks and strata} Local 10 Dafny + 10 Verus from
\texttt{bench/local/q/lab1/STRATUM.toml} (the levers-refused case; the maxauto case is
not drawable without six more Dafny tasks); public 20 DafnyBench hint-free and 20
HumanEval-Verus task forms, with the sets' commits and validation sha256 copied from
\texttt{bench/manifest.toml}; draw seeds 6201 (local), 6202 (DafnyBench), 6203 (HumanEval-Verus); pilot: one sample on each of 80 public
candidates (seeds 90001--90080), a public stratum outside 0.1--0.9 is redrawn once with
the next seed.

\textbf{Generation settings} As the Ollama gate pins them, the same for all arms;
seeds A1 10001--10480, A2 20001--20480, A3 30001--30480, pilot 90001--90080; Claude's
pinned argv.

\textbf{Cost per arm} Tokens in and out, GPU-seconds and wall clock per arm, measured
and reported against the ceiling; the truncation rate per arm as a sensitivity line.

\textbf{Seeds and repeats} Qwen: the seed lists above; Claude: three repeated runs.

\textbf{Cheat surface and gate} My Q09 gate core in the verification container on every
claimed success, the lab gate's lever check, the task forms' spec freeze; grader-only
classes reported \texttt{UNKNOWN}.

\textbf{Exclusions} \texttt{infra-failed}: re-issued with the same id, prompt hash and
seed, at most 3 times, then excluded and counted; \texttt{truncated:<cause>}: a failure
in the primary analysis, excluded in a sensitivity line.

\textbf{Stopping rule} A cap or a window between attempts defers the batch; no outcome
is looked at before the plan's last id is scored; no arm stops on its results.

\textbf{Prediction intervals} (80\%) A1 0.20--0.40, A2 0.25--0.50, A3 0.20--0.45;
A2 $-$ A1 $-$0.05--0.20; A3 $-$ A1 $-$0.10--0.15.

\textbf{Contamination} Local tasks authored 2026-10-07, after both models' release;
the public sets flagged as possibly seen.

\textbf{Disclosure} A small-scale replication with model substitution; Claude is both a
subject and the grader of this plan; one machine, one solver pin; model ids, digests,
argv, prompt hashes and the sandbox manifest recorded per run.

\textbf{Budget arithmetic} Qwen: $60 \text{ tasks} \times 8 \text{ calls} \times 3
\text{ arms} \times 3 \text{ repeats} = 4{,}320$, plus the pilot's $80$, so $4{,}400$
calls in $540 + 80 = 620$ runs; Claude: $20 \times 8 = 160$ calls plus at most 20
sentinels, $\le 180$ calls in 20 runs; equal to the registry's budget. GPU: DV0's single
sample, about 7.3 s per independent call, gives about 9 GPU-hours, a lower bound for A2
and A3; redone per arm from DV4's throughput probe before the runs, with two repeats as
the pre-registered fallback above 10 GPU-hours.

\textbf{MDE} Baseline 0.4 this time (the pilot band's middle), 60 paired tasks, power
0.8, Holm worst case 0.025:
\texttt{stats.mde\_two\_proportions(60, 0.4, alpha=0.025, paired=True)} $= 0.2854$, so
about 29 points; the paired formula is the conservative one (correlated arms discord
less).
\textbf{Batch plan} Prospective file \texttt{research/q12-main.toml}.
This method fixture does not open a batch. Once the policies are final,
record the exact file digest on its own \texttt{plan-sha256:} line and
re-grade the changed pre-registration before executing it. Separately
executed pilot or subset plans each need their own binding and gradebook
row; the registered method stays fixed.
\end{answer}
% PREREG-END

\section*{Problem B2 (lab report)}\label{ans:B2}
\begin{answer}
\textbf{Reproduce:} \verb+<not yet: the runs need DV4 steps 2--4>+
\end{answer}

% ---------------------------------------------------------------------------
\section*{Problem B3 (paper card)}\label{ans:B3}
\begin{answer}
\textbf{Source:} cheatsheet entry 05.10, lines 2148--2171 (the pinned range).

\textbf{Task shape:} 782 verified Dafny programs with their loop invariants and
assertions removed; the model restores hints until Dafny verifies.

\textbf{Oracle:} Dafny verifies with the specification unchanged and no \texttt{assume
false} or \texttt{\{:verify false\}}.

\textbf{What is trusted:} Dafny, the fixed specifications (some weak) and a
two-pattern cheat check.

\textbf{Cheat surface:} Only two cheat patterns are checked; weak specifications; 208
programs verify with no hints at all.

\textbf{Metric:} SR$_{10}$, the share solved within ten attempts with the error fed
back: 67.8\% for Claude 3 Opus at best, 26.9\% when no model is used; first attempts
reach about 54\% and the curve flattens near 65\% by five.

\textbf{Budget:} At most 10 attempts per program, temperature 0.3, 4,096 max tokens.

\textbf{Variance:} The table's $\pm$ values; no repeated runs are reported.

\textbf{Artifact:} The 782-program benchmark, with no training split.

\textbf{Closest arm, and what its budget counts:} A2, attempts per program.
\end{answer}

% ---------------------------------------------------------------------------
\section*{Problem B4 (paper card)}\label{ans:B4}
\begin{answer}
\textbf{Source:} cheatsheet entry 06.1, lines 2312--2338 (the pinned range).

\textbf{Task shape:} Write a whole Isabelle/HOL proof from the theorem statement, or
repair a failed one from Isabelle's error, with a fine-tuned Minerva.

\textbf{Oracle:} Isabelle accepts the proof; \texttt{sorry} and \texttt{oops} are
rejected; 10 s per step.

\textbf{What is trusted:} Isabelle and the theorem-wise PISA split, which favours
copying from nearby proofs.

\textbf{Cheat surface:} \texttt{sorry}/\texttt{oops} are refused; leakage from
pretraining is possible; the best figure is a union with Thor.

\textbf{Metric:} Proof rate, any of $k$ samples accepted: 34.8\% and 40.7\% for 8b at
16 and 64 samples, 47.9\% for 62b with context at 64, 65.7\% for the union with Thor.

\textbf{Budget:} 16 or 64 samples; about 35 TPU-seconds for 64, against 216 s per
proof for step-wise search.

\textbf{Variance:} Not reported, beyond a temperature sweep.

\textbf{Artifact:} None named; the models are closed.

\textbf{Closest arm, and what its budget counts:} A1 for generation, A2 for repair;
samples.
\end{answer}

% ---------------------------------------------------------------------------
\section*{Problem B5 (paper card)}\label{ans:B5}
\begin{answer}
\textbf{Source:} cheatsheet entry 04.8, lines 1681--1706 (the pinned range).

\textbf{Task shape:} Tree search over partial Dafny and Coq programs, one line per
action, with the verifier pruning rejected prefixes; no training.

\textbf{Oracle:} The verifier on complete programs, plus a hidden check lemma that
restates the target property.

\textbf{What is trusted:} Dafny and Coq, and check lemmas written over model-defined
predicates.

\textbf{Cheat surface:} Coarse whole-function checks; Coq prompts carry hints that
nearly give proofs away; the check lemmas rest on the model's own predicates.

\textbf{Metric:} pass@$T$ in sampled tokens; at 5000, about 0.51 Dafny and 0.29 Coq for
VerMCTS against about 0.14 and 0 for whole sampling.

\textbf{Budget:} 5000 sampled tokens per problem; hyperparameters tuned on one of its
own problems.

\textbf{Variance:} Averages read from a figure; no repeats stated.

\textbf{Artifact:} A 15-problem suite, 9 Dafny and 6 Coq.

\textbf{Closest arm, and what its budget counts:} A3; tokens, not calls.
\end{answer}

\end{document}
